\documentclass[10pt,a4paper]{amsart}
\usepackage{amsmath,amssymb,amsthm,graphicx}
% Retained minimal preamble for the Stacks Project chapter coherent.tex (Stacks chapter, pre-delegation retained fragment for stacks-coherent).
% Verbatim excerpt of the chapter preamble (sources/stacks-spaces-descent/preamble.tex): the reference/slogan/history
% environment block (native lines 7-12), Xy-pic setup (native lines 14-19) and the native theorem declarations with
% the mathematical macros (native lines 154-217). Copied mechanically; nothing added, removed or redefined. The
% wrapper class and the amsmath/amssymb/amsthm/graphicx base packages are supplied by the assembler, so the
% class/xr-hyper/hyperref/externaldocument part of the native preamble is omitted.
% For dealing with references we use the comment environment
\usepackage{verbatim}
\newenvironment{reference}{\comment}{\endcomment}
%\newenvironment{reference}{}{}
\newenvironment{slogan}{\comment}{\endcomment}
\newenvironment{history}{\comment}{\endcomment}
% For commutative diagrams we use Xy-pic
\usepackage[all]{xy}

% We use 2cell for 2-commutative diagrams.
\xyoption{2cell}
\UseAllTwocells
% Theorem environments.
%
\theoremstyle{plain}
\newtheorem{theorem}[subsection]{Theorem}
\newtheorem{proposition}[subsection]{Proposition}
\newtheorem{lemma}[subsection]{Lemma}

\theoremstyle{definition}
\newtheorem{definition}[subsection]{Definition}
\newtheorem{example}[subsection]{Example}
\newtheorem{exercise}[subsection]{Exercise}
\newtheorem{situation}[subsection]{Situation}

\theoremstyle{remark}
\newtheorem{remark}[subsection]{Remark}
\newtheorem{remarks}[subsection]{Remarks}

\numberwithin{equation}{subsection}

% Macros
%
\def\lim{\mathop{\mathrm{lim}}\nolimits}
\def\colim{\mathop{\mathrm{colim}}\nolimits}
\def\Spec{\mathop{\mathrm{Spec}}}
\def\Hom{\mathop{\mathrm{Hom}}\nolimits}
\def\Ext{\mathop{\mathrm{Ext}}\nolimits}
\def\SheafHom{\mathop{\mathcal{H}\!\mathit{om}}\nolimits}
\def\SheafExt{\mathop{\mathcal{E}\!\mathit{xt}}\nolimits}
\def\Sch{\mathit{Sch}}
\def\Mor{\mathop{\mathrm{Mor}}\nolimits}
\def\Ob{\mathop{\mathrm{Ob}}\nolimits}
\def\Sh{\mathop{\mathit{Sh}}\nolimits}
\def\NL{\mathop{N\!L}\nolimits}
\def\CH{\mathop{\mathrm{CH}}\nolimits}
\def\proetale{{pro\text{-}\acute{e}tale}}
\def\etale{{\acute{e}tale}}
\def\QCoh{\mathit{QCoh}}
\def\Ker{\mathop{\mathrm{Ker}}}
\def\Im{\mathop{\mathrm{Im}}}
\def\Coker{\mathop{\mathrm{Coker}}}
\def\Coim{\mathop{\mathrm{Coim}}}

% Boxtimes
%
\DeclareMathSymbol{\boxtimes}{\mathbin}{AMSa}{"02}

%
% Macros for moduli stacks/spaces
%
\def\QCohstack{\mathcal{QC}\!\mathit{oh}}
\def\Cohstack{\mathcal{C}\!\mathit{oh}}
\def\Spacesstack{\mathcal{S}\!\mathit{paces}}
\def\Quotfunctor{\mathrm{Quot}}
\def\Hilbfunctor{\mathrm{Hilb}}
\def\Curvesstack{\mathcal{C}\!\mathit{urves}}
\def\Polarizedstack{\mathcal{P}\!\mathit{olarized}}
\def\Complexesstack{\mathcal{C}\!\mathit{omplexes}}
% \Pic is the operator that assigns to X its picard group, usage \Pic(X)
% \Picardstack_{X/B} denotes the Picard stack of X over B
% \Picardfunctor_{X/B} denotes the Picard functor of X over B
\def\Pic{\mathop{\mathrm{Pic}}\nolimits}
\def\Picardstack{\mathcal{P}\!\mathit{ic}}
\def\Picardfunctor{\mathrm{Pic}}
\def\Deformationcategory{\mathcal{D}\!\mathit{ef}}
\usepackage{cleveref}
\crefname{originallemma}{original lemma}{original lemmas}
\Crefname{originallemma}{Original lemma}{Original lemmas}
\makeatletter
\newlabel{lemma-Artin-Rees}{{01YA}{source}{}{sourcepointer-1}{}}
\newlabel{lemma-Artin-Rees@cref}{{[originallemma][0][]01YA}{[1][0][]source}{}{}{}}
\makeatother
\crefname{originallemma}{original lemma}{original lemmas}
\Crefname{originallemma}{Original lemma}{Original lemmas}
\makeatletter
\newlabel{lemma-acc-coherent}{{01Y8}{source}{}{sourcepointer-2}{}}
\newlabel{lemma-acc-coherent@cref}{{[originallemma][0][]01Y8}{[1][0][]source}{}{}{}}
\makeatother
\crefname{originallemma}{original lemma}{original lemmas}
\Crefname{originallemma}{Original lemma}{Original lemmas}
\makeatletter
\newlabel{lemma-coherent-Noetherian}{{01XZ}{source}{}{sourcepointer-3}{}}
\newlabel{lemma-coherent-Noetherian@cref}{{[originallemma][0][]01XZ}{[1][0][]source}{}{}{}}
\makeatother
\crefname{originallemma}{original lemma}{original lemmas}
\Crefname{originallemma}{Original lemma}{Original lemmas}
\makeatletter
\newlabel{modules-lemma-internal-hom}{{01CN}{source}{}{sourcepointer-4}{}}
\newlabel{modules-lemma-internal-hom@cref}{{[originallemma][0][]01CN}{[1][0][]source}{}{}{}}
\makeatother
\crefname{originallemma}{original lemma}{original lemmas}
\Crefname{originallemma}{Original lemma}{Original lemmas}
\makeatletter
\newlabel{properties-lemma-quasi-coherent-colimit-finite-type}{{01PG}{source}{}{sourcepointer-5}{}}
\newlabel{properties-lemma-quasi-coherent-colimit-finite-type@cref}{{[originallemma][0][]01PG}{[1][0][]source}{}{}{}}
\makeatother
\crefname{originallemma}{original lemma}{original lemmas}
\Crefname{originallemma}{Original lemma}{Original lemmas}
\makeatletter
\newlabel{properties-lemma-sections-over-quasi-compact-open}{{01PQ}{source}{}{sourcepointer-6}{}}
\newlabel{properties-lemma-sections-over-quasi-compact-open@cref}{{[originallemma][0][]01PQ}{[1][0][]source}{}{}{}}
\makeatother
\crefname{originallemma}{original lemma}{original lemmas}
\Crefname{originallemma}{Original lemma}{Original lemmas}
\makeatletter
\newlabel{schemes-section-quasi-coherent}{{01LA}{source}{}{sourcepointer-7}{}}
\newlabel{schemes-section-quasi-coherent@cref}{{[originallemma][0][]01LA}{[1][0][]source}{}{}{}}
\makeatother
\title{Hom sheaves over the complement of a closed subscheme: for X a Noetherian scheme, F quasi-coherent, G coherent, I a quasi-coherent ideal sheaf with closed subscheme Z and U = X \textbackslash{} Z, the colimit over n of Hom\_\{O\_X\}(I\textasciicircum{}n G, F) is canonically isomorphic to Hom\_\{O\_U\}(G|\_U, F|\_U); in particular colim\_n Hom\_\{O\_X\}(I\textasciicircum{}n, F) = Gamma(U, F).}
\author{The Stacks Project Authors}
\date{stacks-project:coherent.tex:local-snapshot-20261012; GNU FDL 1.2 or later}
\begin{document}
\maketitle
\noindent\textit{Editorial note.} Statement span is lines 2837-2858 of sources/stacks-coherent/pinned-original-source.tex; proof span is lines 2860-2960, the complete author proof (99 body lines).\par\medskip
\noindent\textit{Editorial note.} Navigation stubs: 7 cited-but-not-included labels are supplied as non-load-bearing pointers with their pinned tags where the snapshot tag index has them.\par\medskip
\noindent\textit{Editorial note.} extra\_packages is empty on purpose: the retained preamble already loads xy with options and a second declaration clashes in the assembler wrapper.\par\medskip
\noindent\textit{Editorial note.} Text is retained verbatim, including the author's own typos and placeholders.\par\medskip
\noindent\textit{Editorial note.} Disclosure per the tightened ruling: the assertion I\textasciicircum{}n K = 0 is marked '(proof omitted)' and injectivity of the general G arrow is omitted; the surjectivity constructions for both the G = O\_X case and the general case are fully written in the span.\par\medskip
\noindent\textit{Editorial note.} Span computed from the \textbackslash{}label marker through the enclosing \textbackslash{}end\{lemma\}; precard\_gate RESULT=OK (probe L9999).\par\medskip
\noindent\textit{Allowed standard prerequisites (editorial).}\begin{itemize}
\item Explicit prerequisite boundary. See why-field; every pointer named here occurs verbatim in the retained statement/proof text.
\end{itemize}\medskip
\ifdefined\hypertarget\hypertarget{sourcepointer-1}{}\fi
\noindent\textit{Source reference.} Original lemma 01YA: sources/stacks-coherent/pinned-original-source.tex, source lines 2802--2815. Cited by the retained statement or proof but not part of the collected spans; supplied only as a navigation pointer. The argument does not depend on its content.\par\medskip
\ifdefined\hypertarget\hypertarget{sourcepointer-2}{}\fi
\noindent\textit{Source reference.} Original lemma 01Y8: sources/stacks-coherent/pinned-original-source.tex, source lines 2766--2776. Cited by the retained statement or proof but not part of the collected spans; supplied only as a navigation pointer. The argument does not depend on its content.\par\medskip
\ifdefined\hypertarget\hypertarget{sourcepointer-3}{}\fi
\noindent\textit{Source reference.} Original lemma 01XZ: sources/stacks-coherent/pinned-original-source.tex, source lines 2442--2455. Cited by the retained statement or proof but not part of the collected spans; supplied only as a navigation pointer. The argument does not depend on its content.\par\medskip
\ifdefined\hypertarget\hypertarget{sourcepointer-4}{}\fi
\noindent\textit{Source reference.} Original lemma 01CN: sources/stacks-coherent/pinned-original-source.tex, source lines 2960--2960. Cited by the retained statement or proof but not part of the collected spans; supplied only as a navigation pointer. The argument does not depend on its content.\par\medskip
\ifdefined\hypertarget\hypertarget{sourcepointer-5}{}\fi
\noindent\textit{Source reference.} Original lemma 01PG: sources/stacks-coherent/pinned-original-source.tex, source lines 2960--2960. Cited by the retained statement or proof but not part of the collected spans; supplied only as a navigation pointer. The argument does not depend on its content.\par\medskip
\ifdefined\hypertarget\hypertarget{sourcepointer-6}{}\fi
\noindent\textit{Source reference.} Original lemma 01PQ: sources/stacks-coherent/pinned-original-source.tex, source lines 2960--2960. Cited by the retained statement or proof but not part of the collected spans; supplied only as a navigation pointer. The argument does not depend on its content.\par\medskip
\ifdefined\hypertarget\hypertarget{sourcepointer-7}{}\fi
\noindent\textit{Source reference.} Original lemma 01LA: sources/stacks-coherent/pinned-original-source.tex, source lines 2960--2960. Cited by the retained statement or proof but not part of the collected spans; supplied only as a navigation pointer. The argument does not depend on its content.\par\medskip
\begin{lemma}
\label{lemma-homs-over-open}
Let $X$ be a Noetherian scheme.
Let $\mathcal{F}$ be a quasi-coherent $\mathcal{O}_X$-module.
Let $\mathcal{G}$ be a coherent $\mathcal{O}_X$-module.
Let $\mathcal{I} \subset \mathcal{O}_X$ be a quasi-coherent sheaf of
ideals. Denote $Z \subset X$ the corresponding closed subscheme and
set $U = X \setminus Z$.
There is a canonical isomorphism
$$
\colim_n \Hom_{\mathcal{O}_X}(\mathcal{I}^n\mathcal{G}, \mathcal{F})
\longrightarrow
\Hom_{\mathcal{O}_U}(\mathcal{G}|_U, \mathcal{F}|_U).
$$
In particular we have an isomorphism
$$
\colim_n \Hom_{\mathcal{O}_X}(
\mathcal{I}^n, \mathcal{F})
\longrightarrow
\Gamma(U, \mathcal{F}).
$$
\end{lemma}

\begin{proof}
We first prove the second map is an isomorphism. It is injective by 
Properties, Lemma \ref{properties-lemma-sections-over-quasi-compact-open}.
Since $\mathcal{F}$ is the union of its coherent submodules, see
Properties, Lemma \ref{properties-lemma-quasi-coherent-colimit-finite-type}
(and Lemma \ref{lemma-coherent-Noetherian})
we may and do assume that $\mathcal{F}$ is coherent to prove surjectivity.
Let $\mathcal{F}_n$ denote the quasi-coherent subsheaf of $\mathcal{F}$
consisting of sections annihilated by $\mathcal{I}^n$,
see Properties, Lemma \ref{properties-lemma-sections-over-quasi-compact-open}.
Since $\mathcal{F}_1 \subset \mathcal{F}_2 \subset \ldots$ we see that
$\mathcal{F}_n = \mathcal{F}_{n + 1} = \ldots $ for some $n \geq 0$
by Lemma \ref{lemma-acc-coherent}. Set $\mathcal{H} = \mathcal{F}_n$
for this $n$. By Artin-Rees (Lemma \ref{lemma-Artin-Rees})
there exists an $c \geq 0$ such that
$\mathcal{I}^m\mathcal{F} \cap \mathcal{H}
\subset \mathcal{I}^{m - c}\mathcal{H}$. Picking $m = n + c$ we get
$\mathcal{I}^m\mathcal{F} \cap \mathcal{H} \subset \mathcal{I}^n\mathcal{H}
= 0$. Thus if we set $\mathcal{F}' = \mathcal{I}^m\mathcal{F}$ then we
see that $\mathcal{F}' \cap \mathcal{F}_n = 0$ and
$\mathcal{F}'|_U = \mathcal{F}|_U$. Note in particular that the subsheaf
$(\mathcal{F}')_N$ of sections annihilated by $\mathcal{I}^N$ is zero
for all $N \geq 0$. Hence by
Properties, Lemma \ref{properties-lemma-sections-over-quasi-compact-open}
we deduce that
the top horizontal arrow in the following commutative
diagram is a bijection:
$$
\xymatrix{
\colim_n \Hom_{\mathcal{O}_X}(
\mathcal{I}^n, \mathcal{F}')
\ar[r] \ar[d] &
\Gamma(U, \mathcal{F}') \ar[d] \\
\colim_n \Hom_{\mathcal{O}_X}(
\mathcal{I}^n, \mathcal{F})
\ar[r] &
\Gamma(U, \mathcal{F})
}
$$
Since also the right vertical arrow is a bijection we conclude that
the bottom horizontal arrow is surjective as desired.

\medskip\noindent
Next, we prove the first arrow of the lemma is a bijection.
By Lemma \ref{lemma-coherent-Noetherian} the sheaf $\mathcal{G}$
is of finite presentation and hence the sheaf
$\mathcal{H} = \SheafHom_{\mathcal{O}_X}(\mathcal{G}, \mathcal{F})$
is quasi-coherent, see
Schemes, Section \ref{schemes-section-quasi-coherent}.
By definition we have
$$
\mathcal{H}(U)
=
\Hom_{\mathcal{O}_U}(\mathcal{G}|_U, \mathcal{F}|_U)
$$
Pick a $\psi$ in the right hand side of the first arrow of the
lemma, i.e.,  $\psi \in \mathcal{H}(U)$. The result just proved applies
to $\mathcal{H}$ and hence there exists an $n \geq 0$ and an
$\varphi : \mathcal{I}^n \to \mathcal{H}$ which recovers
$\psi$ on restriction to $U$. By
Modules, Lemma \ref{modules-lemma-internal-hom}
$\varphi$ corresponds to a map
$$
\varphi :
\mathcal{I}^n \otimes_{\mathcal{O}_X} \mathcal{G}
\longrightarrow
\mathcal{F}.
$$
This is almost what we want except that the source of the arrow
is the tensor product of $\mathcal{I}^n$ and $\mathcal{G}$
and not the product. We will show that, at the cost of increasing $n$,
the difference is irrelevant. Consider the short exact sequence
$$
0 \to \mathcal{K} \to
\mathcal{I}^n \otimes_{\mathcal{O}_X} \mathcal{G} \to
\mathcal{I}^n\mathcal{G} \to 0
$$
where $\mathcal{K}$ is defined as the kernel. Note that
$\mathcal{I}^n\mathcal{K} = 0$ (proof omitted). By Artin-Rees
again we see that
$$
\mathcal{K}
\cap
\mathcal{I}^m(\mathcal{I}^n \otimes_{\mathcal{O}_X} \mathcal{G})
=
0
$$
for some $m$ large enough. In other words we see that
$$
\mathcal{I}^m(\mathcal{I}^n \otimes_{\mathcal{O}_X} \mathcal{G})
\longrightarrow
\mathcal{I}^{n + m}\mathcal{G}
$$
is an isomorphism. Let $\varphi'$ be the restriction of
$\varphi$ to this submodule thought of as a map
$\mathcal{I}^{m + n}\mathcal{G} \to \mathcal{F}$.
Then $\varphi'$ gives an element
of the left hand side of the first arrow of the lemma which
maps to $\psi$ via the arrow. In other words we have proved surjectivity
of the arrow. We omit the proof of injectivity.
\end{proof}

\end{document}
